% Self-contained notation; include this file after \appendix in the main paper.
\section{Identification and stock--flow accounting results}
\label{app:identification-theory}

\subsection{Scope, novelty, and assurance boundary}
\label{app:scope}

This appendix states the mathematical results used to interpret the empirical
design. The rank identities, projection arguments, linear-programming
characterisation, and stock--flow inequalities below are standard consequences
of finite-dimensional linear algebra, convex analysis, and accounting. We do
not claim them as new theorems. The housing-specific contribution is instead
the disciplined application of these results to a prospectively frozen observation
map, the separation of stocks from gross flows, and the use of explicit
response sets and provenance gates before attaching an empirical
interpretation.

Nothing in this appendix supplies a numerical bound for England. In
particular, a response column used in a design audit is not thereby established
as a historical response, a causal impulse response, or a transportable
elasticity. Any empirical interval requires independently sourced response
sets, declared population and time transport assumptions, and a separate
measurement audit. The propositions below describe what follows
\emph{conditional on} their stated matrices and restrictions.

\subsection{Linear response setup}
\label{app:linear-setup}

Let \(y\in\mathbb R^m\) denote an observed vector of moments. The exact linear
response model is
\begin{equation}
  y=f\beta+N\gamma,
  \label{eq:linear-response}
\end{equation}
where \(f\in\mathbb R^m\) is the response column for the focal mechanism,
\(\beta\in\mathbb R\) is its coefficient,
\(N\in\mathbb R^{m\times k}\) collects admitted nuisance and rival response
columns, and \(\gamma\in\mathbb R^k\) contains their coefficients. Zero or
linearly dependent columns of \(N\) are allowed. Unless restrictions are
introduced explicitly, \(\gamma\) is unrestricted.

The phrase \emph{identified against the admitted span} is deliberate. It is a
statement about the columns included in \(N\), not a guarantee that all
economically relevant mechanisms have been included. Enlarging \(N\) can
remove identification; validating the membership of \(N\) is therefore a
substantive part of the empirical design rather than a purely numerical step.

\subsection{A common time path cannot create a new identifying direction}
\label{app:common-path}

Suppose \(a_j\in\mathbb R^q\) is the cross-moment loading for mechanism \(j\),
and all mechanisms are multiplied by the same time path
\(h\in\mathbb R^T\). With moments stacked by date, the corresponding design is
\begin{equation}
  D=h\otimes A,
  \qquad
  A=\begin{bmatrix}a_1&\cdots&a_p\end{bmatrix}
  \in\mathbb R^{q\times p}.
  \label{eq:common-path-design}
\end{equation}

\paragraph{Proposition A.1 (common-path Kronecker rank).}
If \(h\neq0\), then
\begin{equation}
  \operatorname{rank}(h\otimes A)=\operatorname{rank}(A).
  \label{eq:kronecker-rank}
\end{equation}
If \(h=0\), the rank is zero. Consequently, repeating a common ramp or other
common scalar path over time does not separate mechanisms whose cross-moment
loadings were already linearly dependent.

\paragraph{Proof.}
Regard \(h\) as a \(T\times1\) matrix. The rank identity for a Kronecker
product gives
\[
  \operatorname{rank}(h\otimes A)
  =\operatorname{rank}(h)\operatorname{rank}(A).
\]
For nonzero \(h\), \(\operatorname{rank}(h)=1\); for \(h=0\), it is zero.
Equivalently, if \(Ac=0\), then
\((h\otimes A)c=h\otimes(Ac)=0\). Conversely, when \(h\neq0\), select an
index \(t\) with \(h_t\neq0\). The \(t\)-th \(q\)-row block of
\((h\otimes A)c=0\) is \(h_tAc=0\), so \(Ac=0\). Thus the two designs have
the same column null space and, for \(h\neq0\), the same rank.
\(\square\)

\paragraph{Housing-specific interpretation.}
In a housing application, imposing a common temporal ramp on wealth, credit,
supply, or demographic response signatures can add observations without
adding an identifying direction. A longer common path may improve precision
under a valid stochastic model, but it cannot by itself resolve mechanism
collinearity. This conclusion concerns the imposed design and does not assert
that the mechanisms truly share a common path in England.

\subsection{Point identification and its dual certificate}
\label{app:point-identification}

Point identification here means uniqueness of \(\beta\) among all coefficient
pairs satisfying equation~\eqref{eq:linear-response}. It is not a causal claim
and does not address misspecification.

\paragraph{Proposition A.2 (focal identification against an unrestricted nuisance span).}
Assume that equation~\eqref{eq:linear-response} is feasible. The focal
coefficient \(\beta\) is unique for every feasible \(y\) if and only if
\begin{equation}
  f\notin\operatorname{col}(N).
  \label{eq:outside-span}
\end{equation}

\paragraph{Proof.}
Let \((\beta_1,\gamma_1)\) and \((\beta_2,\gamma_2)\) generate the same \(y\).
Their difference satisfies
\begin{equation}
  f(\beta_1-\beta_2)+N(\gamma_1-\gamma_2)=0.
  \label{eq:coefficient-difference}
\end{equation}
If \(f\notin\operatorname{col}(N)\), equation
\eqref{eq:coefficient-difference} cannot hold with
\(\beta_1-\beta_2\neq0\); hence \(\beta_1=\beta_2\).

Conversely, if \(f\in\operatorname{col}(N)\), choose \(a\) such that
\(f=Na\). From any feasible pair \((\beta,\gamma)\), the continuum
\begin{equation}
  (\beta+t,\;\gamma-ta),\qquad t\in\mathbb R,
  \label{eq:observational-equivalence-path}
\end{equation}
generates the same \(y\). Thus \(\beta\) is not unique.
\(\square\)

The unrestricted-nuisance qualification is essential. Sign, box, shape, or
cross-equation restrictions on \(\gamma\) can shrink the feasible values of
\(\beta\), possibly to a singleton, even when \(f\in\operatorname{col}(N)\).
Such a conclusion is identification by the stated restrictions and must be
reported as such.

\paragraph{Proposition A.3 (dual discriminator equivalence).}
The following statements are equivalent:
\begin{enumerate}
  \item \(f\notin\operatorname{col}(N)\);
  \item there exists \(d\in\mathbb R^m\) such that
        \(d^{\mathsf T}N=0\) and \(d^{\mathsf T}f=1\).
\end{enumerate}
When they hold, \(\beta=d^{\mathsf T}y\) in the exact model.

\paragraph{Proof.}
Let \(P_N\) be the Euclidean orthogonal projector onto
\(\operatorname{col}(N)\), and define \(r=(I-P_N)f\). If
\(f\notin\operatorname{col}(N)\), then \(r\neq0\),
\(N^{\mathsf T}r=0\), and
\[
  r^{\mathsf T}f
  =r^{\mathsf T}(P_Nf+r)
  =r^{\mathsf T}r>0.
\]
Therefore \(d=r/(r^{\mathsf T}r)\) has the required properties. Conversely,
if \(f=Na\) for some \(a\), then every \(d\) satisfying
\(d^{\mathsf T}N=0\) also satisfies \(d^{\mathsf T}f=d^{\mathsf T}Na=0\),
so no unit-loading discriminator exists. Finally, multiplying
equation~\eqref{eq:linear-response} by \(d^{\mathsf T}\) yields
\(d^{\mathsf T}y=\beta\).
\(\square\)

The vector \(d\) is a constructive diagnostic: it identifies the contrast of
observed moments that annihilates every admitted nuisance column. It is not a
substantive validation of those columns. A discriminator can disappear when a
new, credible rival response is added to \(N\).

\subsection{Observation maps and missing-coordinate loss}
\label{app:observation-map}

Let \(y^*\in\mathbb R^r\) be a latent or more richly measured response,
\begin{equation}
  y^*=f^*\beta+N^*\gamma,
  \label{eq:latent-response}
\end{equation}
and let the available observations be \(y=My^*\) for a known linear map
\(M\in\mathbb R^{m\times r}\). The observed response columns are then
\(Mf^*\) and \(MN^*\).

\paragraph{Proposition A.4 (identification after an observation map).}
Under unrestricted nuisance coefficients, the focal coefficient is identified
from the observed vector if and only if
\begin{equation}
  Mf^*\notin\operatorname{col}(MN^*).
  \label{eq:observed-identification}
\end{equation}
Equivalently, observed identification fails if and only if there exists
\(a\in\mathbb R^k\) such that
\begin{equation}
  f^*-N^*a\in\ker(M).
  \label{eq:kernel-equivalence}
\end{equation}
Moreover,
\(\operatorname{rank}(M[f^*\;N^*])
\leq\operatorname{rank}([f^*\;N^*])\): an observation map can preserve or
destroy design rank, but cannot increase the rank of the latent design.

\paragraph{Proof.}
Apply Proposition~A.2 to the observed equation
\[
  y=(Mf^*)\beta+(MN^*)\gamma.
\]
Failure of condition~\eqref{eq:observed-identification} means that
\(Mf^*=MN^*a\) for some \(a\), which is exactly
\(M(f^*-N^*a)=0\), proving equation~\eqref{eq:kernel-equivalence}. The rank
inequality is the standard inequality
\(\operatorname{rank}(MD)\leq\operatorname{rank}(D)\), applied to
\(D=[f^*\;N^*]\).
\(\square\)

If \(M\) selects observed coordinates, \(\ker(M)\) consists of vectors
supported on missing coordinates. Proposition~A.4 then makes the loss
mechanism explicit: a focal-minus-rival contrast may be visible in the richer
system but live entirely in coordinates absent from the public observation
map. Observing more dates of the same mapped coordinates need not recover that
contrast.

\paragraph{Housing-specific interpretation.}
Candidate missing bridges include a mapping from property transitions to
household tenure, a mapping from ownership or letting states into measured
prices, and a mapping from lifecycle groups into the published age categories.
The proposition says only that such coordinates \emph{could} distinguish
mechanisms. It does not establish that a particular restricted-access linkage
is accurate, licence-compatible, or sufficient, nor that it would distinguish
wealth from every admitted rival.

\subsection{Weighted least-squares profile geometry}
\label{app:weighted-profile}

For noisy or overidentified moments, let \(W\in\mathbb R^{m\times m}\) be
symmetric positive definite and consider
\begin{equation}
  Q(\beta,\gamma)
  =(y-f\beta-N\gamma)^{\mathsf T}
    W(y-f\beta-N\gamma).
  \label{eq:weighted-objective}
\end{equation}
Let \(C\) be any nonsingular square root satisfying \(C^{\mathsf T}C=W\),
write \(\widetilde y=Cy\), \(\widetilde f=Cf\), and
\(\widetilde N=CN\), and let \(P_{\widetilde N}\) be the Euclidean
orthogonal projector onto \(\operatorname{col}(\widetilde N)\). Define
\begin{equation}
  r_y=(I-P_{\widetilde N})\widetilde y,
  \qquad
  r_f=(I-P_{\widetilde N})\widetilde f.
  \label{eq:weighted-residuals}
\end{equation}

\paragraph{Proposition A.5 (profile objective and curvature).}
With \(\gamma\) unrestricted,
\begin{equation}
  Q_p(\beta)
  :=\inf_{\gamma}Q(\beta,\gamma)
  =\lVert r_y-\beta r_f\rVert_2^2.
  \label{eq:profile-objective}
\end{equation}
If \(r_f\neq0\), then
\begin{align}
  \widehat\beta
    &=\frac{r_f^{\mathsf T}r_y}{r_f^{\mathsf T}r_f},
    \label{eq:profile-beta}\\
  Q_p(\beta)
    &=Q_p(\widehat\beta)
      +(r_f^{\mathsf T}r_f)(\beta-\widehat\beta)^2.
    \label{eq:profile-quadratic}
\end{align}
If \(r_f=0\), the profile is flat:
\(Q_p(\beta)=\lVert r_y\rVert_2^2\) for all \(\beta\). Furthermore,
\(r_f=0\) if and only if \(f\in\operatorname{col}(N)\).

\paragraph{Proof.}
For fixed \(\beta\), minimizing equation~\eqref{eq:weighted-objective} over
\(\gamma\) is Euclidean least squares after premultiplication by \(C\). Its
residual is the projection of \(\widetilde y-\beta\widetilde f\) onto the
orthogonal complement of \(\operatorname{col}(\widetilde N)\), giving
equation~\eqref{eq:profile-objective}. Expanding the square and completing it
in \(\beta\) gives equations~\eqref{eq:profile-beta} and
\eqref{eq:profile-quadratic} when \(r_f\neq0\). When \(r_f=0\), the term
depending on \(\beta\) vanishes. Finally, because \(C\) is nonsingular,
\(Cf\in\operatorname{col}(CN)\) if and only if
\(f\in\operatorname{col}(N)\).
\(\square\)

Thus \(r_f^{\mathsf T}r_f\) is the profile curvature, or design information,
for the focal coefficient against the admitted span. Near-zero curvature
warns that apparently precise coefficient estimates may be dominated by
scaling, regularisation, or restrictions. The formula is an algebraic
property of the supplied \(f,N,W\); it is not evidence that \(W\) is the true
sampling precision or that \(\widehat\beta\) is causal.

\subsection{Sharp conditional focal sets}
\label{app:sharp-set}

Now impose explicit bounds. Let
\begin{equation}
  \Gamma=\{\gamma\in\mathbb R^k:\ell\leq\gamma\leq u\},
  \qquad
  B=[b_L,b_U],
  \label{eq:box-sets}
\end{equation}
where inequalities are componentwise and either focal endpoint may be
infinite. The conditional feasible set and its focal projection are
\begin{align}
  \Theta(y)
    &=\{(\beta,\gamma):
          y=f\beta+N\gamma,\ \beta\in B,\ \gamma\in\Gamma\},
    \label{eq:conditional-feasible-set}\\
  \mathcal I_\beta(y)
    &=\{\beta:\text{there exists }\gamma
          \text{ with }(\beta,\gamma)\in\Theta(y)\}.
    \label{eq:focal-identified-set}
\end{align}

\paragraph{Proposition A.6 (sharp set and linear-program endpoints).}
The set \(\mathcal I_\beta(y)\) is an interval, a ray, the real line, a
singleton, or the empty set. When nonempty, its lower and upper endpoints in
the extended real line are the values of the linear programs
\begin{align}
  \underline\beta
    &=\inf_{\beta,\gamma}\ \beta
      &&\text{subject to }(\beta,\gamma)\in\Theta(y),
    \label{eq:lower-lp}\\
  \overline\beta
    &=\sup_{\beta,\gamma}\ \beta
      &&\text{subject to }(\beta,\gamma)\in\Theta(y).
    \label{eq:upper-lp}
\end{align}
Every finite optimum is attained. These endpoints are sharp conditional on
the equality, the columns, and the bounds: every value between two feasible
finite endpoint values is generated by a feasible coefficient pair.

\paragraph{Proof.}
The set \(\Theta(y)\) is a polyhedron because it is defined by linear
equalities and inequalities. Its projection onto the scalar \(\beta\)-axis is
also a polyhedron, hence has one of the listed forms. Linear objectives over a
nonempty polyhedron attain any finite optimum, giving
equations~\eqref{eq:lower-lp}--\eqref{eq:upper-lp}. If
\((\beta_L,\gamma_L)\) and \((\beta_U,\gamma_U)\) are feasible endpoint
witnesses, then for any \(\lambda\in[0,1]\),
\[
  \lambda(\beta_L,\gamma_L)
  +(1-\lambda)(\beta_U,\gamma_U)
\]
is feasible by convexity. Its focal coordinate covers every value between
\(\beta_L\) and \(\beta_U\), proving sharpness.
\(\square\)

An empty set is not a failed optimisation to be silently repaired: it rejects
the joint equality and restrictions. An unbounded endpoint is likewise a
scientific result. Most importantly, mathematical sharpness is conditional
sharpness. Boxes chosen after inspecting \(y\), response columns calibrated
from the focal column, or rival sets omitted for convenience do not become
empirically credible merely because the endpoint programs are solved exactly.

\subsection{Why stocks do not identify gross tenure flows}
\label{app:stock-flow}

Let \(x_{jt}\geq0\) be the level of the stock in tenure state \(j\) at the
start of period \(t\), for \(j=1,\ldots,K\). Let \(F_{ij,t}\geq0\) be the
number of units transitioning internally from state \(i\) to state \(j\),
\(i\neq j\), and let \(q_{jt}\) collect net external additions to state \(j\)
(entries less exits, new supply less removals, and any other declared boundary
terms). The accounting identity is
\begin{equation}
  d_{jt}
  :=x_{j,t+1}-x_{jt}-q_{jt}
  =\sum_{i\neq j}F_{ij,t}-\sum_{k\neq j}F_{jk,t}.
  \label{eq:stock-flow-accounting}
\end{equation}
When all internal states and external boundary terms are consistently
measured, \(\sum_jd_{jt}=0\).

\paragraph{Proposition A.7 (gross-flow non-identification from stocks).}
Suppose \(K\geq2\), only the net changes \(d_t\) are observed, and no capacity,
turnover, or upper-bound restriction is imposed on \(F_t\). If one nonnegative
flow matrix satisfies equation~\eqref{eq:stock-flow-accounting}, infinitely
many do. In particular, for any distinct \(a,b\) and any \(c\geq0\), adding
\(c\) to both \(F_{ab,t}\) and \(F_{ba,t}\) leaves every \(d_{jt}\)
unchanged. Hence neither total gross internal flow nor the two opposing flows
of a pair are point-identified by stock changes alone.

\paragraph{Proof.}
The added \(a\)-to-\(b\) flow subtracts \(c\) from the balance of state \(a\)
and adds \(c\) to the balance of state \(b\). The added reverse flow does the
opposite. The net effect is zero in every state, while total gross flow rises
by \(2c\). Since \(c\) is arbitrary, the observed stock changes admit a
continuum of observationally equivalent gross flows.
\(\square\)

\paragraph{Proposition A.8 (sharp lower bound on total internal flow).}
Assume that \(q_t\) is known, \(\sum_jd_{jt}=0\), and transitions are allowed
between every pair of distinct states. Let
\begin{equation}
  G_t=\sum_i\sum_{j\neq i}F_{ij,t}
  \label{eq:total-gross-flow}
\end{equation}
count each internal transition once. Then
\begin{equation}
  G_t\geq
  \sum_j(d_{jt})_+
  =\frac12\sum_j\lvert d_{jt}\rvert,
  \label{eq:gross-flow-lower-bound}
\end{equation}
and the bound is sharp. State-specific gross inflow and outflow satisfy
\begin{equation}
  \sum_{i\neq j}F_{ij,t}\geq(d_{jt})_+,
  \qquad
  \sum_{k\neq j}F_{jk,t}\geq(-d_{jt})_+.
  \label{eq:state-flow-lower-bounds}
\end{equation}

\paragraph{Proof.}
Write \(I_j=\sum_{i\neq j}F_{ij,t}\) and
\(O_j=\sum_{k\neq j}F_{jk,t}\). Since \(d_j=I_j-O_j\) and both terms are
nonnegative, \(I_j\geq(d_j)_+\) and \(O_j\geq(-d_j)_+\). Summing the first
inequality over \(j\) gives \(G_t=\sum_jI_j\geq\sum_j(d_j)_+\). The zero-sum
condition implies
\(\sum_j(d_j)_+=\sum_j(-d_j)_+=\frac12\sum_j|d_j|\).

To attain the bound, treat every state with \(d_j<0\) as supplying \(-d_j\)
units and every state with \(d_j>0\) as demanding \(d_j\) units. Because total
supply equals total demand and the transition graph is complete, a
nonnegative transportation plan can send exactly the required amount from
shrinking to growing states, with no counterflow or cycle. Its total flow is
\(\sum_j(d_j)_+\).
\(\square\)

If the transition graph is restricted, equation~\eqref{eq:gross-flow-lower-bound}
remains a lower bound when a feasible flow exists, but it need not be sharp.
For two states in a closed system,
\(d_2=F_{12}-F_{21}\), so
\(F_{12}\geq(d_2)_+\) and \(F_{21}\geq(-d_2)_+\); even there, the two gross
flows remain unbounded above without additional restrictions. With three or
more states, a rise in one state and fall in another generally does not impose
a positive lower bound on the \emph{direct} flow between that particular pair,
because the balance can be routed through other states.

Two further qualifications matter in housing data. First, if \(q_t\) is
unobserved and unrestricted, internal flow need not have a positive lower bound:
setting \(q_t=x_{t+1}-x_t\) and \(F_t=0\) satisfies the accounting identity.
If instead \(q_t\) belongs to a declared feasible set \(\mathcal Q_t\), the
best stock-based lower bound is
\begin{equation}
  \inf_{q\in\mathcal Q_t:\,\boldsymbol 1^{\mathsf T}
         (x_{t+1}-x_t-q)=0}
  \frac12\lVert x_{t+1}-x_t-q\rVert_1,
  \label{eq:external-flow-robust-bound}
\end{equation}
provided the associated internal-flow problem is feasible. Second, changes in
tenure \emph{shares} alone do not determine changes in tenure levels when the
total number of households or dwellings changes. The level denominator and
the household--dwelling measurement bridge must be specified before applying
the accounting bounds.

\paragraph{Housing-specific interpretation.}
An owner-occupation stock decline and a private-rental stock increase are net
accounting facts, not measurements of gross owner-to-landlord purchases,
landlord-to-owner sales, or household moves. Applying
equation~\eqref{eq:gross-flow-lower-bound} requires compatible levels, boundary
terms, units, geography, and timing. It cannot attribute the residual flow to
the wealth of buyers without an additional buyer-resource measurement and a
credible behavioural design.

\subsection{A cautious design value for an additional measurement}
\label{app:measurement-value}

The preceding results can rank candidate measurements without claiming that a
dataset has policy or welfare value. Return to the latent columns \(f^*,N^*\).
For an observation map \(M\) and positive-definite weight \(W\), define the
focal separation index
\begin{equation}
  J(M,W)
  =\min_{a\in\mathbb R^k}
    \bigl\|W^{1/2}M(f^*-N^*a)\bigr\|_2^2.
  \label{eq:separation-index}
\end{equation}
This is the squared weighted distance from \(Mf^*\) to
\(\operatorname{col}(MN^*)\). It is positive exactly when the focal column is
identified against the admitted nuisance span under \(M\).

Let \((M_0,W_0)\) describe existing moments. A candidate additional
measurement has map \(H\) and positive-definite weight \(W_H\), with block
weights used only as a design convention. Define
\begin{align}
  J_0
    &=\min_a
      \bigl\|W_0^{1/2}M_0(f^*-N^*a)\bigr\|_2^2,
    \label{eq:old-separation}\\
  J_1
    &=\min_a\left\{
      \bigl\|W_0^{1/2}M_0(f^*-N^*a)\bigr\|_2^2
      +\bigl\|W_H^{1/2}H(f^*-N^*a)\bigr\|_2^2
      \right\}.
    \label{eq:new-separation}
\end{align}

\paragraph{Proposition A.9 (nonnegative design value of an added coordinate).}
Under the fixed latent columns, fixed nuisance set, and block-weight convention
above, \(J_1\geq J_0\). Equality holds if and only if at least one minimiser of
the old problem also satisfies
\begin{equation}
  H(f^*-N^*a)=0.
  \label{eq:no-new-separation}
\end{equation}
In particular, if \(J_0=0\), the additional measurement restores point
identification against the admitted span if and only if there is no
\(a\) satisfying both
\begin{equation}
  M_0(f^*-N^*a)=0
  \quad\text{and}\quad
  H(f^*-N^*a)=0.
  \label{eq:break-all-equivalences}
\end{equation}

\paragraph{Proof.}
For each \(a\), the objective in equation~\eqref{eq:new-separation} is the old
nonnegative objective plus a nonnegative term. Taking minima therefore gives
\(J_1\geq J_0\). Both are finite-dimensional positive-semidefinite quadratic
problems and attain their minima. Equality requires a minimiser of the new
problem to attain the old minimum while making the added nonnegative term zero;
this is exactly condition~\eqref{eq:no-new-separation}. When \(J_0=0\), the
old minimisers are precisely the solutions to the first equality in
equation~\eqref{eq:break-all-equivalences}. The stacked design has positive
separation exactly when none of them also satisfies the second equality.
\(\square\)

Proposition~A.9 defines \emph{design value}, not the expected social value of
information. It assumes the old and new response columns are correctly
specified on a common scale, the nuisance set remains fixed, and the new
measurement does not introduce selection, linkage, or transport error. A
practical acquisition decision must also account for coverage, sampling error,
revision, linkage quality, licence constraints, privacy, cost, and the risk that
a defensible enlargement of \(N^*\) absorbs the apparent new separation.

In the housing application, a candidate measurement is most informative for
mechanism separation when it observes a coordinate on which every currently
equivalent focal--rival contrast differs. Buyer-resource links, gross property
status changes, or household--dwelling reconciliation may be candidates, but
Proposition~A.9 does not presume that any available dataset measures them well
or that any one of them is sufficient. The proposition supports a
value-of-information \emph{design audit}; it does not support an empirical
England contribution estimate.

\subsection{Implications for the empirical claim ceiling}
\label{app:claim-ceiling}

The results imply a hierarchy of defensible statements:
\begin{enumerate}
  \item A full-rank maintained design establishes only that its maintained
        columns are algebraically distinct.
  \item If a credible rival span contains the focal column, unrestricted
        unique attribution fails, even when the maintained design alone has
        full rank.
  \item A nonzero dual discriminator certifies separation only against the
        currently admitted columns and observation map.
  \item Box or shape restrictions may deliver a sharp conditional interval,
        but the interval is empirical only after the restrictions and response
        sets are independently sourced and transported.
  \item Stock changes can support accounting lower bounds under declared
        boundary assumptions; they do not reveal gross flows or buyer wealth.
  \item Additional data have positive mechanism-separation value only to the
        extent that they break the relevant observational equivalences after
        measurement error and an enlarged rival set are considered.
\end{enumerate}

Accordingly, the formal results justify a pivot from unrestricted unique
attribution toward partial identification and measurement design. They do not
reverse a failed identification gate, establish that a particular mechanism
caused England's tenure changes, or authorise a numerical historical
contribution bound without the separate empirical response-set package
described in Section~\ref{app:scope}.
